\documentclass[11pt]{article}
\usepackage[margin=1in]{geometry}
\usepackage{times}
\title{SE333 Final Project — Reflection}
\author{Your Name}
\date{\today}
\begin{document}
\maketitle
\begin{abstract}
This short reflection summarizes the goals, methodology, results, and lessons learned while building an MCP-based testing agent for the SE333 final project. Keep this under two pages when compiled.
\end{abstract}

\section{Introduction}
Briefly state the project goals: build an MCP server that generates, runs, and iterates unit tests to improve JaCoCo coverage for a Maven project.

\section{Methodology}
Describe the main components: the FastMCP server (`server.py`), Java test generation (`generate_basic_junit`), coverage analysis (`parse_jacoco`), and automation (git tools, CI).

\section{Results}
During development I ran the test suite and generated JaCoCo reports. Current results from `sample-maven/target/site/jacoco/jacoco.xml` (generated locally):

\begin{itemize}
  \item Line coverage: 100.00% (2/2 lines covered)
  \item Branch coverage: N/A (no branch counters present)
\end{itemize}

Continuous Integration: the repository includes a GitHub Actions workflow (`.github/workflows/maven.yml`) which runs on push and PR using JDK 21 and executes `mvn -f sample-maven/pom.xml test`.

\section{Lessons Learned}
Short bullets about AI-assisted testing, limitations of naive test generation, and integration pitfalls.

\section{Future Work}
Mention planned extensions (spec-based test generation, AI code review agent, dashboard).

\section{Appendix}
Provide build instructions to reproduce the PDF and coverage report.
\vspace{6pt}
\textbf{How to build PDF:}
\begin{verbatim}
# On a system with LaTeX installed:
pdflatex -interaction=nonstopmode -output-directory=docs docs/reflection.tex
# Output: docs/reflection.pdf
\end{verbatim}

\end{document}
